% TOP_PROBLEM: {"id": 3181, "title": "Matchings extend to Hamiltonian cycles in hypercubes", "category_id": 3, "category": {"id": 3, "name": "graph_theory", "display_name": "Graph Theory", "description": "Problems involving graphs, networks, and their properties.", "slug": "graph-theory", "order_index": 3, "created_at": "2026-07-31T15:26:25.671Z"}, "status": "open", "problem_number": "OPG-600", "set_id": 12, "statusQualification": "The intended hypercube is Q_d with d at least two, made explicit following the primary formulation."}
% FIRST_POSTED: 2026-10-01
% ENTRY_KIND: statement_only
% UnsolvedMath ID 3181; source catalogue code OPG-600
% !TeX program = lualatex
% Definition and sources only; this file is not an LLM solution attempt.
\documentclass[11pt,a4paper]{article}
\usepackage[margin=25mm]{geometry}
\usepackage{fontspec}
\setmainfont{lmroman10-regular.otf}[BoldFont=lmroman10-bold.otf,
 ItalicFont=lmroman10-italic.otf,BoldItalicFont=lmroman10-bolditalic.otf]
\usepackage{amsmath,amssymb,mathtools}
\usepackage{unicode-math}
\setmathfont{latinmodern-math.otf}
\AtBeginDocument{\let\setminus\smallsetminus}
\usepackage{xurl,hyperref}
\hypersetup{colorlinks=true,urlcolor=blue}
\setcounter{secnumdepth}{0}
\setlength{\parindent}{0pt}
\setlength{\parskip}{5pt}
\setlength{\emergencystretch}{3em}
\begin{document}
\section*{Matchings extend to Hamiltonian cycles in hypercubes}\label{3181}
{\small UnsolvedMath ID 3181 · OPG-600 · Graph Theory · OpenGarden}
\subsection{Definitions and mathematical statement}
For an integer $d\geq2$, the $d$-dimensional hypercube
$Q_d$ is the simple graph with vertex set $\{0,1\}^d$.
Two bit vectors are adjacent exactly when they differ
in one coordinate. Thus $Q_d$ has $2^d$ vertices and
degree $d$ at every vertex. A matching is a set
$M\subseteq E(Q_d)$ whose edges have pairwise disjoint
endpoint sets. The empty matching is allowed, and a
matching need not cover all vertices.

A Hamilton cycle is a simple cycle passing through
every vertex of $Q_d$ exactly once. To extend the
matching means to find such a cycle $C$ containing
each prescribed matching edge:
\[
 M\subseteq E(C).
\]
The conjecture asks whether this is possible for every
integer $d\geq2$ and every matching $M$ of $Q_d$.

Equivalently, the task is to cyclically order all binary
vectors so consecutive vectors differ in one coordinate,
while requiring the endpoints of every edge of $M$ to
appear consecutively in that cyclic order. Prescribed
edges may point in different coordinate directions and
may have any total size up to $2^{d-1}$. No choice of
orientation of the matching is imposed. The lower bound
on $d$ is necessary under the ordinary simple-cycle
convention, since $Q_1$ is a single edge. A theorem for
perfect matchings alone would not directly address all
matchings required here.
\subsection{Short English statement}
Can every collection of pairwise disjoint hypercube edges be included in one cycle visiting all hypercube vertices?
\subsection{Sources}
\begin{itemize}
\item[S1] ULAM.AI and the UnsolvedMath Contributors, supplied problems.json, record 3181 (OPG-600), OpenGarden collection. Catalogue identity and source transcription; CC BY 4.0.
\url{https://huggingface.co/datasets/ulamai/UnsolvedMath}.
\item[S2] Open Problem Garden, Matchings extend to Hamiltonian cycles in hypercubes. Original problem and discussion; the mathematical formulation below is expanded from the supplied source record.
\url{http://www.openproblemgarden.org/op/matchings_extends_to_hamilton_cycles_in_hypercubes}.
\item[S3] J. Fink and T. Mütze, Matchings in hypercubes extend to long cycles, arXiv:2401.01769, revised 2025; precise Ruskey–Savage formulation.
\url{https://arxiv.org/abs/2401.01769}.
\end{itemize}
\end{document}
